\section{Retractions from the branches}\label{branches}
Let $V_i$ ($1\le i\le s$) be DVRs, with maximal ideals $\mathfrak n_i=(\pi_i)$ and identified residue field $k$. Let $R=\prod_k V_i$ be the subring of tuples with common residue in $k$. Its maximal ideal is $J=\prod_i\mathfrak n_i$. The coordinate maps are onto, each summand of $J$ is Noetherian over its branch, and consequently $R$ is Noetherian. Its normalization is $\prod_i V_i$, since the quotient by $R$ is a finite-dimensional $k$-space and inverting $(\pi_i)_i$ gives the product of the fraction fields. Thus $R$ is reduced of dimension one, and its minimal primes are the coordinate kernels $P_i$.
\begin{theorem}[Branch criterion]\label{branch-split}
The ring $R=\prod_k V_i$ has $\K$ in every characteristic. When the branches are complete, a module-finite extension $R\subseteq S$ splits exactly when some local factor of $S$ is nonzero at every minimal prime of $R$.
\end{theorem}
\begin{lemma}\label{glue}
Every module-finite extension $R\subseteq T$ with $T$ local admits a retraction.
\end{lemma}
\begin{proof}
Put $T_i=T/P_iT$. Then $V_i\to T_i$ is injective. Indeed the tuple supported on branch $i$ with value $\pi_i$ annihilates $P_iT$. If $r\in R\cap P_iT$, multiplication by that tuple gives $\pi_i r_i=0$ in $V_i$. Thus $r\in P_i$. The $V_i$-torsion in $T_i$ is a proper ideal $N_i$, hence is contained in the unique maximal ideal of $T_i$. The closed fibers $T_i/\mathfrak n_iT_i$ all identify with $E=T/JT$.

Let $L$ be the residue field of $T$ and choose a $k$-linear functional $L\to k$ taking $1$ to $1$. The resulting functional $\lambda:E\to k$ kills the image of each $N_i$. Since $T_i/N_i$ is finite free over $V_i$, the induced functional on its closed fiber lifts, by choosing a basis, to a $V_i$-linear functional. Its value at $1$ is congruent to $1$ modulo $\mathfrak n_i$. Divide by this unit and compose with $T\to T_i/N_i$ to get $\psi_i:T\to V_i$. All the $\psi_i$ reduce to $\lambda$ and take $1$ to $1$. Thus $(\psi_i)_i$ takes values in $R$ and is the desired $R$-linear retraction.
\end{proof}
\begin{proof}[Proof of Theorem~\ref{branch-split}]
If the $V_i$ are complete then $R$ is complete, since $J^n=\prod_i\mathfrak n_i^n$. A finite $R$-algebra $S$ is therefore a product of nonzero local algebras $T_j$. A factor is nonzero on every generic branch exactly when the map $R\to T_j$ has zero kernel. Indeed, the generic localizations are fields and $R$ embeds in their product. Such a factor gives a retraction by Lemma~\ref{glue}. Conversely, decomposing a retraction as a sum of functionals on the factors gives $1=\sum_j\rho_j(1_{T_j})$. Some summand is a unit in the local ring $R$, forcing the corresponding factor map to be injective.

If $S$ has finite projective dimension, so does every $T_j$. The alternating sum of the ranks in a finite free resolution of $T_j$ equals its vector-space dimension on each generic branch. Some factor has positive rank on the first branch since $R\subseteq S$. It therefore has positive rank on every branch and yields the required retraction.

In general $\widehat R=\prod_k\widehat V_i$ by the same computation of $J^n$. Flat base change preserves the extension and its finite free resolution. For any flat extension $R\to R'$ the evaluation ideal extends to the evaluation ideal over $R'$, since $S$ is finitely presented and $\Hom_R(S,R)$ commutes with flat base change. Faithful flatness of completion thus descends the retraction.
\end{proof}

\subsection{Separable residue fields and intrinsic curves}

Let $(R,\m,k)$ be an excellent one-dimensional local ring and put
$A=(\widehat R)_{\mathrm{red}}$.  Its normalization is finite and has
the form $\overline A=\prod_iV_i$, where the $V_i$ are complete DVRs.
Their residue fields $k_i$ are finite over $k$.  We use the condition
\[
 \tag{SB}\label{separable-branches}
 \text{Every residue field extension }k_i/k\text{ is separable.}
\]
The use of the reduced completion makes this condition meaningful
before reducedness or seminormality of $R$ is known.  It is automatic
when $k$ is perfect, but the argument only needs
\eqref{separable-branches}.

We first recall the classical conductor description of a complete
seminormal curve \cite{Swan,EH}.

\begin{lemma}\label{conductor-description}
Let $R$ be a complete one-dimensional seminormal local ring with finite
normalization $V=\prod_iV_i$, where the $V_i$ are complete DVRs.  Put
$J=\prod_i\mathfrak n_i$ and $k_i=V_i/\mathfrak n_i$.  Then $J=\m$
as subsets of $V$, and $R$ is the inverse image of the diagonal image
of $k$ under the residue map $V\to\prod_i k_i$.
\end{lemma}
\begin{proof}
The finite module $V/R$ is supported at the closed point.  Its
annihilator, the conductor $I=(R:V)$, therefore contains a power of
$\m$.  Since $I$ is also an ideal of $V$ and $\sqrt{\m V}=J$, it
contains a power of $J$.

For $b\in J$, all sufficiently high powers of $b$ lie in $I\subseteq R$.
The square--cube criterion for seminormality \cite{Swan} says that
$c^2,c^3\in R$, for $c\in Q(R)$, imply $c\in R$.  If all powers $b^n$
with $n\geq N$ belong to $R$, applying this criterion to $b^n$ gives
the same conclusion for $n\geq\lceil N/2\rceil$.  Iteration yields
$b\in R$, so $J\subseteq R$.  The quotient $R/J$ is a reduced Artinian
local subring of $\prod_i k_i$, hence a field.  It is therefore the
residue field $k$ of $R$, proving both assertions.
\end{proof}

\begin{proof}[Proof of Theorem B]
Suppose first that $R$ is seminormal and satisfies
\eqref{separable-branches}.  Completion preserves seminormality for
excellent local rings \cite[Cor.~5.3]{GT}.  The evaluation-ideal
argument in the proof of Theorem~\ref{branch-split} permits descent, so we may assume
that $R$ is complete.  Let $V=\prod_iV_i$ be its normalization.
Lemma~\ref{conductor-description} identifies $R$ with the inverse image
of $k\subseteq\prod_i k_i$.

By \eqref{separable-branches}, a finite Galois extension $l/k$ splits
all the fields $k_i$.  Lift the monic minimal polynomial of a primitive
element of $l/k$ to a monic polynomial $f\in R[t]$.  Then
$R'=R[t]/(f)$ is finite free, local and complete, with residue field
$l$.  Its derivative is a unit at the class of $t$, as can be checked
in $l$, so $R'/R$ is finite \'etale.

Put $V'=V\otimes_RR'$.  On each branch, the field factors of
\[
 k_i\otimes_k l\simeq l^{[k_i:k]}
\]
lift by henselianity to local factors of the finite \'etale
$V_i$-algebra $V_i\otimes_RR'$.  Each factor is a complete DVR because it is
flat over $V_i$, its quotient by $\pi_i$ is $l$, and its maximal ideal
is generated by the regular element $\pi_i$.  Tensoring the exact
sequence
\[
 0\longrightarrow R\longrightarrow V\oplus k
 \xrightarrow{(v,c)\mapsto\overline v-c}\prod_i k_i
 \longrightarrow0
\]
with the flat algebra $R'$ identifies $R'$ with the fiber product of
these DVR factors over their common residue field $l$.  Theorem~\ref{branch-split}
gives $\K$ for $R'$, and faithful flatness descends it to $R$.

For necessity in residue characteristic two, suppose that $R$ is
Cohen--Macaulay, satisfies $\K$, and that two is either zero or a
nonzerodivisor.  If $R$ were not weakly normal, Lemma~\ref{wn}(3) would
give $q\in Q(R)\setminus R$ with $2q,q^2\in R$.  Theorem~\ref{realize} would
produce a nonsplit module-finite extension of finite projective dimension,
a contradiction.  Thus $R$ is weakly normal and in particular
seminormal.  Together with the sufficient direction this proves the
asserted equivalence under \eqref{separable-branches}.  The statement
about $F$-purity follows from the next corollary.
\end{proof}

The corresponding Frobenius consequence holds in every positive equal
characteristic.  It fits the classical structure theory of
one-dimensional $F$-pure rings \cite{GW,Gib}.

\begin{corollary}\label{separable-fpure}
Let $R$ be an excellent one-dimensional Cohen--Macaulay local ring of
characteristic $p>0$ satisfying \eqref{separable-branches}.  Then
seminormality, weak normality and $F$-purity are equivalent.
No finiteness hypothesis on Frobenius is required.
\end{corollary}
\begin{proof}
Weak normality implies seminormality.  If Frobenius is pure and
$q=b/a\in Q(R)$ satisfies $q^p\in R$, with $a$ a nonzerodivisor, then
$b^p\in a^pR$.  Tensoring Frobenius with $R/aR$ gives an injection
\[
 R/aR\longrightarrow F_*(R/a^pR),\qquad
 b+aR\longmapsto F_*(b^p+a^pR).
\]
Hence $b\in aR$, so $q\in R$, and Lemma~\ref{wn} gives weak normality.

Conversely, suppose that $R$ is seminormal.  The completion and finite
\'etale extension used above give a faithfully flat local extension
$R\to R'$ such that $R'$ is complete and is a fiber product of DVRs
with common residue field $l$.  A coefficient field of $R'$ maps to
a coefficient field of each branch.  Choosing uniformizers therefore
identifies
\[
 R'\simeq\prod_l l[[t_i]].
\]
The field $l$ need not be perfect.  Regard $F_*l$ as an $l$-vector
space with scalar action $a\cdot F_*c=F_*(a^pc)$.  Since $F_*1\ne0$,
choose an $l$-linear functional
\[
 \lambda:F_*l\longrightarrow l,\qquad \lambda(F_*1)=1.
\]
On each branch define
\[
 C_i\!\left(F_*\sum_{n\geq0}c_nt_i^n\right)
   =\sum_{m\geq0}\lambda(F_*c_{pm})t_i^m.
\]
Additivity and the identity
$\lambda(F_*(a^pc))=a\lambda(F_*c)$ show coefficientwise that $C_i$
is $l[[t_i]]$-linear for the Frobenius module structure.  It takes
$F_*1$ to $1$ and hence splits Frobenius.  The maps $C_i$ induce the
same functional $\lambda$ on the common residue field, so their
product restricts to a Frobenius splitting $F_*R'\to R'$.

Finally, $R\to R'\to F_*R'$ is pure: the first map is faithfully flat
and the second is split.  This composite factors through the
Frobenius map $R\to F_*R$.  After tensoring with any $R$-module,
injectivity of the composite forces injectivity of the first factor.
Thus $R$ is $F$-pure.
\end{proof}

For instance $k[[x_1,\ldots,x_s]]/(x_ix_j:i<j)$ identifies with $\prod_k k[[x_i]]$ by restriction to the axes. Theorem~\ref{branch-split} gives $\K$ for all $s$. For $s\ge2$ the normalization itself is a nonsplit module-finite extension because each local factor misses the other branches. Its projective dimension is infinite, since a single branch has generic ranks $1$ and $0$ at different minimal primes and hence cannot admit a finite free resolution.
